\fichatitulo{46 · AVALIAÇÃO JURÍDICA}{Workshop NLLP · 2024}{Groundedness em QA jurídico}
{\small Trautmann et al.. \textit{Groundedness em QA jurídico}. Workshop NLLP · 2024. \href{https://aclanthology.org/2024.nllp-1.14/}{Fonte consultada}.\par}

\campo{Problema e objetivo}
Como avaliar se respostas jurídicas permanecem sustentadas pelo contexto recuperado?

\campo{Principal contribuição · síntese interpretativa}
Compara métricas de similaridade, modelos NLI e estratégias de prompting de LLM para classificar groundedness.

\campo{Método / natureza da evidência}
Benchmark de perguntas e respostas jurídicas com corpus de classificação de grounding; mede também latência.

\campo{Resultado ou argumento principal}
O melhor método alcança macro-F1 de 0,80 nas condições do estudo; o resultado é específico ao corpus, às classes e aos modelos testados.

\campo{Excertos-chave · seleção literal}
\begin{itemize}
\excerto{1}{The best method achieves a macro-F1 score of 0.8}{Resumo.}
\end{itemize}

\campo{Limitações e análise crítica}
A avaliação jurídica é uma ponte metodológica relevante, mas o domínio e o inglês limitam a transferência direta. A dissertação deverá construir exemplos e critérios próprios para discursos parlamentares.

\campo{Aproveitamento na dissertação · proposta}
Metodologia: fundamentar uma rubrica de fidelidade ao contexto e comparar juiz automático com avaliadores humanos.

\registro{Fonte consultada nesta preparação: PDF publicado; consulta focal do resumo, método, resultados e conclusão. Chave: \texttt{trautmannEtAl2024groundedness}.}
